\chapter{Matériel programmable II~: conceptions de trading}\label{ch:nw:programmable-hardware-ii-trading-designs}

La conception de ce chapitre réagit deux cycles d'horloge après que le dernier octet d'un prix l'a atteinte~: \qty{12.8}{\nano\second} à 156.25~MHz,
pour le premier ordre comme pour le dix-millième, que le marché soit calme ou en pleine rafale. Elle n'a pas de cache à manquer ni de thread à réveiller,
et l'ordre qu'elle envoie part d'un \omterm{def:nw:programmable-hardware-ii-trading-designs:template}{gabarit d'ordre} dans lequel trois champs sont complétés. Elle ne peut pas non plus faire grand-chose~: elle compare un prix à un
seuil, pour un seul instrument, et tout ce qu'elle sait du monde a été écrit dans ses registres par le logiciel un instant plus tôt.
Cette combinaison (très rapide, très étroite, pilotée par le logiciel) est ce que les firmes entendent par un \omterm{def:nw:programmable-hardware-ii-trading-designs:hybrid}{système de trading hybride}.

Ce chapitre construit un \omterm{def:nw:programmable-hardware-ii-trading-designs:trigger}{déclencheur matériel} sur le flux du simulateur de bourse, avec ses contrôles de risque et son \omterm{def:nw:programmable-hardware-ii-trading-designs:template}{gabarit d'ordre}, en SystemVerilog
sur la boîte à outils du chapitre~6, et le vérifie de quatre façons~: dans Verilator, dans Icarus, face à un modèle de cycles en C++20 et face à un modèle Python,
ordre par ordre et cycle par cycle. Il place ensuite le déclencheur dans le chemin wire-to-wire de la firme et se demande ce qu'il apporte.

\section{Décodage du flux en matériel}

Un gestionnaire logiciel (livre~13, chapitre~18) reçoit un datagramme entier, puis parcourt ses messages. Un décodeur matériel parcourt les
messages au fur et à mesure que les octets arrivent, huit par cycle, en ne conservant que l'état dont il a besoin entre les cycles~: sa position dans l'en-tête du paquet, la
longueur du message courant, sa progression à l'intérieur de celui-ci, et les champs qu'il a déjà capturés. Dans le flux du simulateur, comme dans le
format de flux de la bourse qu'il suit, un message d'ajout d'ordre fait 36 octets, son prix occupant les quatre derniers~; le décodeur connaît donc tout
d'un ordre dès le cycle où son dernier octet arrive, alors que le reste du paquet (autres messages, séquence de contrôle de la trame) peut encore
rester à venir.

\omcode{code/firm/hwtrade/hdl/hwt_trigger.sv}{47}{75}{Logique combinatoire du décodeur~: une machine à états sérielle par octet, déroulée sur les
huit voies du beat. Elle capture le type, l'instrument, le sens, la taille et le prix, et, au dernier octet du message, décide s'il faut
déclencher.}\label{lst:nw:programmable-hardware-ii-trading-designs:parse}

\begin{definition}[Déclencheur matériel, décision cut-through]\label{def:nw:programmable-hardware-ii-trading-designs:trigger}
Un \emph{déclencheur matériel}\index{déclencheur matériel} est une logique placée sur une \omterm{def:nw:networking-for-trading:frame}{carte réseau} qui surveille un flux de données de marché pour une condition fixée
à l'avance par le logiciel (un cours franchissant un seuil, une transaction sur un instrument donné) et, lorsqu'elle se produit, lance un ordre sans
l'intervention de l'hôte. Une \emph{décision cut-through}\index{décision cut-through} est une décision prise dès que les octets dont elle dépend
sont arrivés, avant que la trame qui les porte ne soit achevée ou vérifiée.
\end{definition}

La machine à états déroulée du \cref{lst:nw:programmable-hardware-ii-trading-designs:parse} est la conception correcte la plus simple, et elle est
lente en matériel~: huit étapes d'octet enchaînées en un seul cycle forment un long chemin critique. Des décodeurs plus rapides exploitent le format~: les messages commencent à
des décalages connus par rapport à leurs préfixes de longueur, si bien qu'une logique dédiée peut repérer un début de message dans chaque voie d'un coup, et les messages à format fixe
peuvent être reconnus par comparaison de mots entiers. C'est le travail de la \omterm{def:nw:programmable-hardware-i-architecture-and-toolchain:timing}{fermeture temporelle} du chapitre~6, et c'est pourquoi les décodeurs de production sont
propres au format d'une seule plateforme.

\section{Construction du carnet d'ordres en matériel}

Un déclencheur sur des messages isolés n'a pas besoin de carnet. Une stratégie qui négocie sur le meilleur cours acheteur ou le meilleur cours vendeur en a besoin, et un carnet en matériel est une
structure de données différente de celles du livre~13, chapitre~19. Les ordres identifiés par leur référence exigent une table de hachage en \omterm{def:nw:programmable-hardware-i-architecture-and-toolchain:bram}{block RAM} ou en mémoire externe
mémoire~; les niveaux de prix exigent un tableau indexé par le décalage du prix par rapport à une référence, le meilleur niveau étant trouvé par un codeur de priorité sur les
niveaux occupés~; chaque mise à jour doit s'achever dans le délai entre deux messages, sinon la conception doit les mettre en file d'attente. Les firmes qui construisent des carnets
en matériel ne conservent que quelques niveaux près du meilleur prix, pour les rares instruments qui en ont besoin, et laissent le reste au logiciel.

\section{Contrôles de risque en cycles}

Le matériel d'une firme envoie des ordres pour le compte d'un broker-dealer, qui doit garder les contrôles préalables à la négociation «~sous le contrôle direct et exclusif
~» du courtier (livre~11, chapitre~27), des contrôles qui rejettent les ordres dépassant des «~paramètres de prix ou de taille appropriés,
au cas par cas ou sur une courte période de temps~». En matériel, ces contrôles se réduisent à des comparaisons et des compteurs, et ils coûtent une
étape du pipeline.

\begin{definition}[Carte des registres]\label{def:nw:programmable-hardware-ii-trading-designs:regmap}
Une \emph{carte des registres}\index{carte des registres} est l'ensemble des adresses auxquelles une conception matérielle expose sa configuration et son état au
logiciel (seuils, limites, bits d'activation et de kill, compteurs), lus et écrits sur le bus de l'hôte, et le document qui les
définit.
\end{definition}

\omcode{code/firm/hwtrade/hdl/hwt_trigger.sv}{108}{135}{L'étape de risque et le gabarit d'ordre~: une limite de taille, un bit de kill et un seau à
jetons, puis l'identifiant, la taille et le prix de l'ordre insérés.}\label{lst:nw:programmable-hardware-ii-trading-designs:risk}

\begin{proposition}[Rafale d'un seau à jetons]\label{prop:nw:programmable-hardware-ii-trading-designs:bucket}
Un seau à jetons de capacité $B$, rempli d'un jeton tous les $R$ cycles, laisse sortir au plus $B + \lfloor t/R\rfloor$ ordres sur tout intervalle
de $t$ cycles. Si le logiciel a besoin de $t$ cycles pour repérer un emballement et activer le bit de kill, le matériel peut envoyer au plus ce nombre d'ordres
avant de s'arrêter, et aucun après.
\end{proposition}

\begin{proof}
Le seau contient au plus $B$ jetons au début de l'intervalle et en gagne au plus un tous les $R$ cycles~; chaque ordre en consomme un. Le
bit de kill masque l'étape de sortie dans le cycle où il est activé.
\end{proof}

La proposition est le document d'exigences d'une couche de risque matérielle. À 156.25~MHz, avec le seau de 4 jetons du chapitre, rempli
tous les 64 cycles, un moniteur logiciel qui réagit en \qty{20}{\micro\second} laisse sortir au plus 52 ordres~; celui qui réagit en
\qty{100}{\micro\second} en laisse sortir 248. Les limites qui comptent en matériel sont celles qui s'imposent avant que le logiciel puisse agir.

\section{Génération d'ordres à partir de gabarits}

\begin{definition}[Gabarit d'ordre]\label{def:nw:programmable-hardware-ii-trading-designs:template}
Un \emph{gabarit d'ordre}\index{gabarit d'ordre} (order template) est un message d'ordre préparé à l'avance, complet à l'exception des champs qui dépendent de
du déclencheur (identifiant client de l'ordre, taille, prix), que le matériel insère au fur et à mesure de l'envoi.
\end{definition}

\begin{definition}[Moteur de déchargement TCP]\label{def:nw:programmable-hardware-ii-trading-designs:toe}
Un \emph{moteur de déchargement TCP}\index{moteur de déchargement TCP} (TCP offload engine) implante le protocole TCP (numéros de séquence, accusés de réception, retransmission)
dans le matériel d'une \omterm{def:nw:networking-for-trading:frame}{carte réseau}, de sorte que la carte puisse émettre et recevoir sur une connexion TCP sans passer par l'hôte.
\end{definition}

La saisie d'ordres passe par TCP (chapitre~1), si bien qu'une carte qui envoie seule des ordres doit gérer l'état de la connexion~: son prochain
numéro de séquence, ses accusés de réception, ses retransmissions. Un \omterm{def:nw:programmable-hardware-ii-trading-designs:toe}{moteur de déchargement TCP} s'en charge~; le \omterm{def:nw:programmable-hardware-ii-trading-designs:template}{gabarit d'ordre} transporte alors les en-têtes TCP et
IP, avec numéro de séquence et sommes de contrôle complétés au moment de l'envoi. Le logiciel qui gère la session (ouverture de session, heartbeat, la
couche de session de la plateforme) partage la connexion avec le matériel par le même moteur.

\begin{dated}{2026-09}{Un chiffre publié de tick-to-trade matériel}\label{dat:nw:programmable-hardware-ii-trading-designs:stac}
Dans un benchmark STAC-T0 publié en juin 2024, un système FPGA construit sur une carte AMD Alveo UL3524 dotée d'un cœur TCP et UDP (une solution d'Exegy et AMD
) a été mesuré à un minimum de \qty{13.9}{\nano\second} pour des trames de 507 octets et \qty{14.1}{\nano\second} pour des trames de 68 octets, depuis
«~le dernier bit de données entrantes nécessaire à une décision de trading jusqu'au premier bit de l'ordre sortant simulé~». La définition
exclut le temps d'arrivée des bits entrants, qu'aucune conception, celle du chapitre comprise, ne peut éviter.
\end{dated}

\section{Conceptions hybrides~: ce qui reste dans le logiciel et pourquoi}

\begin{definition}[Système de trading hybride]\label{def:nw:programmable-hardware-ii-trading-designs:hybrid}
Un \emph{système de trading hybride}\index{système de trading hybride} répartit une stratégie entre le logiciel, qui calcule son modèle, ses positions
et ses paramètres à son propre rythme, et le matériel, qui applique à flux tendu des décisions précalculées aux données de marché au \omterm{def:nw:networking-for-trading:ser}{débit de ligne}~; le logiciel écrit dans les
registres du matériel, et en lit les compteurs et les copies de ce qu'il a envoyé.
\end{definition}

La répartition découle de ce que chaque côté fait bien. Le matériel est rapide et figé~: chaque fonctionnalité coûte de la logique, chaque modification une construction de
quelques heures et une vérification de plusieurs jours, et il ne peut pas retenir beaucoup d'état. Le logiciel est lent et souple~: il peut modéliser, apprendre, compenser les positions
entre plateformes et changer en quelques minutes. Le logiciel décide donc de \emph{ce} qui vaudrait la peine d'être fait («~acheter jusqu'à 700 à 100.02 ou mieux si
quelqu'un propose~»), et le matériel décide de \emph{quand} (le cycle où le cours vendeur arrive). Tout ce qui demande un historique, de nombreux instruments ou
un jugement reste dans le logiciel~; la seule comparaison qui doit d'abord avoir lieu passe sur la carte.

\begin{omfigure}
\begin{tikzpicture}[font=\footnotesize,b/.style={draw,rounded corners=2pt,minimum height=0.7cm,align=center,inner sep=3pt},>=Stealth]
\node[b,fill=gray!10] (feed) at (0,0) {flux (UDP)};
\node[b,fill=omDef!18] (dec) at (2.5,0) {décodage\\ (étape 1)};
\node[b,fill=omThm!15] (risk) at (5.4,0) {risque~: taille,\\ seau, kill\\ (étape 2)};
\node[b,fill=omMeth!20] (tpl) at (8.3,0) {gabarit\\ d'ordre\\ (étape 3)};
\node[b,fill=gray!10] (out) at (11.0,0) {ordres\\ (TCP offload)};
\node[b,fill=omDef!8,minimum width=5.6cm] (sw) at (4.0,-2.4) {logiciel sur l'hôte~: modèle, positions, limites};
\node[b,fill=gray!8] (regs) at (4.0,-1.25) {carte des registres~: seuils, limites, kill, compteurs};
\draw[->,thick] (feed) -- (dec);
\draw[->,thick] (dec) -- node[above,font=\scriptsize]{signal} (risk);
\draw[->,thick] (risk) -- (tpl);
\draw[->,thick] (tpl) -- (out);
\draw[->,omThm] (sw) -- (regs);
\draw[->,omThm] (regs.north -| risk.south) -- (risk.south);
\draw[->,omThm] (regs.north -| dec.south) -- (dec.south);
\draw[->,gray,dashed] (tpl.south) |- node[pos=0.25,right,font=\scriptsize]{copies des ordres} (sw.east);
\end{tikzpicture}

\omcaption{La conception hybride du chapitre. Trois étapes du pipeline sur la carte décodent le flux, vérifient chaque ordre candidat et l'insèrent
dans un gabarit~; le logiciel sur l'hôte écrit les seuils, les limites et le bit de kill par la \omterm{def:nw:programmable-hardware-ii-trading-designs:regmap}{carte des registres} et reçoit les copies de tout
ce que la carte envoie.}\label{fig:nw:programmable-hardware-ii-trading-designs:hybrid}
\end{omfigure}

\begin{omfigure}
\begin{tikzpicture}[font=\scriptsize,c/.style={draw,minimum height=0.55cm,minimum width=1.05cm,align=center,inner sep=1pt}]
\foreach \i in {0,...,9} {\node at (1.1*\i,1.25) {\i};}
\node[anchor=east] at (-0.7,1.25) {cycle};
\foreach \i in {0,...,7} {\node[c,fill=omDef!15] at (1.1*\i,0.55) {mot \i};}
\node[anchor=east] at (-0.7,0.55) {entrée};
\node[c,fill=omDef!35] at (1.1*7,-0.15) {prix};
\node[anchor=east] at (-0.7,-0.15) {décodage};
\node[c,fill=omThm!25] at (1.1*8,-0.85) {risque};
\node[anchor=east] at (-0.7,-0.85) {risque};
\node[c,fill=omMeth!30] at (1.1*9,-1.55) {ordre};
\node[anchor=east] at (-0.7,-1.55) {gabarit};
\end{tikzpicture}

\omcaption{Un paquet contenant un message d'ajout d'ordre (58~octets~: 20 d'en-tête, 2 de longueur, 36 de message) traverse la conception, à raison d'un mot de
8~octets par cycle. Le dernier octet du cours arrive dans le mot~7~; le déclencheur est enregistré à la fin de ce cycle, la décision de risque un cycle
plus tard, et l'ordre au cycle suivant~: deux cycles après le cours, dix après le premier octet du paquet.}\label{fig:nw:programmable-hardware-ii-trading-designs:cycles}
\end{omfigure}

\section{Tutoriel~: un déclencheur sur le flux du simulateur, de quatre façons}\label{tut:nw:programmable-hardware-ii-trading-designs:tut}

\begin{tutorial}
\textbf{Objectif.} Faire tourner le déclencheur sur deux secondes du flux du simulateur de bourse, dans Verilator, dans Icarus, dans un modèle à cycles en C++20 et dans
un modèle Python, vérifier que les quatre envoient les mêmes ordres aux mêmes cycles, et placer la conception sur le chemin wire-to-wire.
\textbf{État final~:} \cref{tab:nw:programmable-hardware-ii-trading-designs:runs}, \cref{fig:nw:programmable-hardware-ii-trading-designs:race}
et les tests au vert de \texttt{firm.hwtrade}.
\begin{tutsteps}
\item \textbf{Le flux.} \texttt{make\_hwtrade\_fixture.py} fait tourner le simulateur du livre~10 pendant deux secondes depuis l'ouverture (992 paquets)
  et les enregistre~; \texttt{beats\_of} les découpe en mots avec deux cycles d'inactivité entre les paquets (la journée est compressée~: le seau compte
  des cycles, et non des secondes).
\item \textbf{Quatre modèles.} \texttt{build\_verilator}, \texttt{build\_icarus}, \texttt{CycleModel} et \texttt{cpp/firm\_hwtrade.hpp}
  (\cref{lst:nw:programmable-hardware-ii-trading-designs:model}) exécutent les mêmes réglages~; les tests comparent leurs sorties ligne à ligne.
\item \textbf{La référence.} \texttt{triggers()} analyse les mêmes paquets message par message en Python~: chaque message qualifié doit être
  soit un ordre, soit un rejet.
\item \textbf{Le chemin.} \texttt{nw\_hwtrade.paths()} construit la variante matérielle de \texttt{firm.wirepath} (la cage du chapitre~2, les
  émetteurs-récepteurs de la carte, trois cycles) à côté du chemin logiciel du chapitre~5, et \texttt{win\_probability} confronte chacun à un concurrent~;
  \texttt{python fig\_hwtrade.py} écrit les données du graphique.
\end{tutsteps}
\textbf{Pour aller plus loin.} Régler le seau à un jeton rechargé toutes les 10\,000 cycles et compter les ordres~; déplacer le bit de kill au
cycle précédant le dixième ordre et vérifier que exactement neuf partent.
\end{tutorial}

\omcode{code/firm/hwtrade/cpp/firm_hwtrade.hpp}{20}{44}{Front d'horloge du modèle à cycles en C++20~: étape 3, puis étape 2, puis étape 1, chacune
lisant les registres qu'a écrits le front précédent, comme le fait le matériel.}\label{lst:nw:programmable-hardware-ii-trading-designs:model}

\begin{table}[ht]
\centering\small
\begin{tabular}{rrrrr}
\toprule
Seuil (prix $\times 10^4$) & Taille maximale & Kill positionné au cycle & Ordres & Rejets \\
\midrule
1\,000\,200 & 10\,000 & jamais & 108 & 5 \\
1\,000\,300 & 10\,000 & jamais & 133 & 27 \\
1\,000\,300 & 600 & jamais & 131 & 29 \\
1\,000\,300 & 10\,000 & 3\,000 & 41 & 119 \\
\bottomrule
\end{tabular}
\caption{Le déclencheur sur deux secondes du flux du simulateur (instrument~1~; ajouts d'ordres sell side au seuil ou en dessous), avec un seau
de 4 jetons rechargés toutes les 64 cycles. Verilator, Icarus, le modèle à cycles en C++20 et le modèle Python donnent les mêmes ordres aux mêmes
cycles~; ordres et rejets totalisent les messages qualifiés que trouve la référence Python. Données~:
le jeu d'essai de \texttt{firm.hwtrade}.}\label{tab:nw:programmable-hardware-ii-trading-designs:runs}
\end{table}

Le tableau montre chaque contrôle à l'œuvre. Un seuil plus bas déclenche moins souvent~; le seau rejette la plupart des déclenchements supplémentaires du
seuil plus haut, car ils arrivent par rafales~; la limite de taille en rejette deux de plus~; et le bit de kill, positionné au cycle 3\,000, arrête tous les ordres suivants. Sur
le chemin wire-to-wire du chapitre~5 (\cref{fig:nw:programmable-hardware-ii-trading-designs:race}), le chemin matériel se compose de la fibre de la cage,
des émetteurs-récepteurs de la carte (\qty{120}{\nano\second} dans chaque sens, une valeur de modèle) et de trois cycles~: \qty{0.82}{\micro\second} à chaque quantile.
Le chemin logiciel a une médiane de \qty{4.4}{\micro\second} et un 99\up{e}~centile de \qty{35}{\micro\second}. Face à un concurrent dont la latence a une
médiane de \qty{1}{\micro\second} (une loi log-normale d'écart-type 0.3, un modèle), le chemin matériel gagne 75~\% des courses, et le chemin logiciel
une sur dix mille.

\begin{omfigure}
\begin{tikzpicture}
\begin{groupplot}[group style={group size=2 by 1,horizontal sep=1.8cm},width=6.6cm,height=5.0cm,tick label style={font=\small},
  label style={font=\small},grid=major,grid style={gray!25},table/col sep=comma]
\nextgroupplot[xmin=0,xmax=1,ymode=log,ymin=0.5,ymax=200,xlabel={quantile},ylabel={temps bout en bout (\unit{\micro\second}, log)},
  ytick={1,10,100},yticklabels={1,10,100},legend columns=2,legend style={font=\footnotesize,at={(1.15,-0.32)},anchor=north,
  draw=none,fill=none,/tikz/every even column/.append style={column sep=8pt}}]
\addplot[omThm,very thick,mark=*,mark size=1.3pt] table[x=q,y=software_us]{figdata/networks/07-programmable-hardware-ii-trading-designs/paths.csv};
\addplot[omDef,very thick,mark=square*,mark size=1.3pt] table[x=q,y=hardware_us]{figdata/networks/07-programmable-hardware-ii-trading-designs/paths.csv};
\legend{chemin logiciel (chapitre~5),chemin matériel}
\nextgroupplot[xmode=log,xmin=250,xmax=12000,ymin=0,ymax=1.05,xlabel={latence du concurrent, médiane (ns, log)},ylabel={probabilité de gagner},
  xtick={300,1000,3000,10000},xticklabels={300,1\,000,3\,000,10\,000}]
\addplot[omThm,very thick,mark=*,mark size=1.3pt] table[x=median_ns,y=software]{figdata/networks/07-programmable-hardware-ii-trading-designs/race.csv};
\addplot[omDef,very thick,mark=square*,mark size=1.3pt] table[x=median_ns,y=hardware]{figdata/networks/07-programmable-hardware-ii-trading-designs/race.csv};
\end{groupplot}
\end{tikzpicture}

\omcaption{Simulation~: latence wire-to-wire du chemin logiciel du chapitre~5 (étapes mesurées du livre~13, la carte et la cage) et du
chemin matériel (la cage, les transceivers de la carte à \qty{120}{\nano\second} dans chaque sens, une valeur de modèle, et trois cycles à 156.25~MHz). À gauche~:
quantiles~; le chemin matériel est une constante. À droite~: la probabilité de battre un concurrent dont la latence est lognormale, de
médiane donnée et une dispersion de 0.3 (un modèle). Données~: \texttt{fig\_hwtrade.py}.}\label{fig:nw:programmable-hardware-ii-trading-designs:race}
\end{omfigure}

\section{Construction~: le déclencheur matériel}\label{bld:nw:programmable-hardware-ii-trading-designs:hwtrade}

\begin{build}
\bfield{Rôle} Un \omterm{def:nw:programmable-hardware-ii-trading-designs:trigger}{déclencheur matériel} complet et testé~: décodage du flux de données, contrôles de risque et \omterm{def:nw:programmable-hardware-ii-trading-designs:template}{gabarit d'ordre} (order template), avec les trois modèles qui le rendent
fiable, et une variante matérielle du chemin wire-to-wire de la firme pour le plan du chapitre~29.
\bfield{Interface} \texttt{hdl/hwt\_trigger.sv} (paramètres \texttt{LOC}, \texttt{BURST}, \texttt{REFILL}~; registres \texttt{thresh},
\texttt{max\_qty}, \texttt{kill}~; sorties~: identifiant, taille et prix de l'ordre, et compteur de rejets)~; \omterm{def:nw:programmable-hardware-i-architecture-and-toolchain:tb}{bancs de test} (testbench) \texttt{tb/hwt\_tb.cpp} et
\texttt{tb/hwt\_tb.sv}~; C++20 \texttt{firm::hwtrade::Model}~; Python \texttt{firm\_hwtrade}~: \texttt{beats\_of}, \texttt{write\_stim},
\texttt{CycleModel}, \texttt{triggers}, \texttt{build\_verilator}, \texttt{build\_icarus}, \texttt{run\_verilator}, \texttt{run\_icarus},
\texttt{parse\_output}, \texttt{wirepath\_stage}.
\bfield{Règles} Au plus un ordre par message qualifiant~; rien ne sort tant que le bit kill est positionné~; les tailles au-delà de la limite sont rejetées et
non tronquées~; le seau à jetons ne prête jamais~; les modèles concordent avec le registre HDL, registre par registre.
\bfield{Tests d'acceptation} \texttt{code/firm/hwtrade/}~: le modèle Python reproduit les quatre réglages du jeu de référence, et chaque message qualifiant
est un ordre ou un rejet~; Verilator et Icarus impriment les mêmes ordres que le modèle Python~; le modèle C++20 reproduit le jeu de référence~;
la limite de taille à sa frontière~; rien après le bit kill~; le rationnement du seau~; deux cycles entre le prix et l'ordre.
\bfield{Pour aller plus loin} Maintenir le meilleur cours vendeur de l'instrument à partir des messages add, execute et delete, et déclencher dessus~; regrouper les trois contrôles en
un seul cycle et mesurer le chemin critique dont il aurait besoin.
\end{build}

\begin{omsources}
\item STAC, «~STAC-T0~», résultats obtenus avec une solution FPGA d'Exegy et AMD (2024).
\item 17 CFR 240.15c3-5 (la Market Access Rule).
\item Nasdaq, \emph{TotalView-ITCH 5.0}, spécification (structure Add Order)~; le \path{code/firm/exchsim/PROTOCOL.md} du simulateur de bourse.
\end{omsources}

\section{Exercices}

\begin{exercise}[$\star$]\label{exo:nw:programmable-hardware-ii-trading-designs:1}
Un paquet transporte un ordre add de 36~octets derrière l'en-tête de 20~octets et le champ de longueur de 2~octets. À quel battement se trouve le dernier octet du prix, et combien de
nanosecondes après le premier battement du paquet le circuit émet-il l'ordre à 156.25~MHz~?
\end{exercise}

\begin{exercise}[$\star$]\label{exo:nw:programmable-hardware-ii-trading-designs:2}
Avec un seau de 4 jetons rechargés toutes les 64 cycles à 156.25~MHz, quel est le débit maximal soutenu d'ordres~?
\end{exercise}

\begin{exercise}[$\star$]\label{exo:nw:programmable-hardware-ii-trading-designs:3}
Combien d'ordres le circuit du chapitre peut-il envoyer avant qu'un moniteur logiciel, qui réagit en \qty{50}{\micro\second}, ne positionne le bit kill~?
\end{exercise}

\begin{exercise}[$\star\star$]\label{exo:nw:programmable-hardware-ii-trading-designs:4}
Dans le \cref{tab:nw:programmable-hardware-ii-trading-designs:runs}, pourquoi relever le seuil de 1\,000\,200 à 1\,000\,300 ajoute-t-il 25 ordres mais
22 rejets~?
\end{exercise}

\begin{exercise}[$\star\star$]\label{exo:nw:programmable-hardware-ii-trading-designs:5}
Pourquoi le circuit rejette-t-il les ordres dont la taille dépasse la limite, au lieu de les envoyer à cette limite~?
\end{exercise}

\begin{exercise}[$\star\star$]\label{exo:nw:programmable-hardware-ii-trading-designs:6}
Citer trois décisions que la stratégie du chapitre laisse au logiciel, et pourquoi chacune serait difficile à réaliser dans le matériel.
\end{exercise}

\begin{exercise}[$\star\star\star$]\label{exo:nw:programmable-hardware-ii-trading-designs:7}
\emph{Programmation.} Avec \texttt{firm.hwtrade.CycleModel}, fixer le seau à 1 jeton rechargé toutes les 10\,000 cycles et compter les ordres au seuil
1\,000\,300. Puis positionner le bit kill un cycle avant le dixième ordre de la simulation sans restriction, et recompter.
\end{exercise}

\begin{exercise}[$\star\star\star$]\label{exo:nw:programmable-hardware-ii-trading-designs:8}
\emph{Trouver la faille.} «~Notre matériel déclenche en 13 nanosecondes, donc notre latence wire-to-wire est de 13 nanosecondes.~»
\end{exercise}

\section{Problème~: Le déclencheur qui ne doit pas se déclencher deux fois}

\begin{problem}[{Problème du week-end --- un chemin hybride, une course et un emballement}]\label{pb:nw:programmable-hardware-ii-trading-designs:1}
Une firme transfère la décision d'une stratégie sensible à la latence sur le \omterm{def:nw:programmable-hardware-ii-trading-designs:trigger}{déclencheur matériel} du chapitre. Utiliser les chiffres du chapitre~: la
médiane wire-to-wire du chemin logiciel \qty{4.39}{\micro\second} et le 99\up{e}~centile \qty{35.3}{\micro\second}~; le chemin matériel \qty{0.82}{\micro\second}
(\qty{263}{\nano\second} de cage à l'entrée, \qty{120}{\nano\second} de carte à l'entrée, trois cycles à 156.25~MHz, \qty{120}{\nano\second} de carte en sortie,
\qty{295}{\nano\second} de cage en sortie)~; un concurrent dont la latence est lognormale~; un seau de 4 jetons rechargé toutes les 64 cycles.

\medskip\noindent\textbf{Partie I --- Le chemin.}
\begin{enumerate}
\item Additionner les étapes du chemin matériel.
\item Quelle part en représente la décision elle-même~?
\item Quel est le 99\up{e}~centile du chemin matériel, et pourquoi~?
\item De combien le chemin matériel devance-t-il le chemin logiciel à la médiane, et au 99\up{e}~centile~?
\end{enumerate}

\medskip\noindent\textbf{Partie II --- La course.}
\begin{enumerate}[resume]
\item Face à un concurrent de médiane \qty{1}{\micro\second}, quelles sont les probabilités de gagner des deux chemins (celles de la
  simulation du chapitre)~?
\item Et face à un concurrent de médiane \qty{5}{\micro\second}~?
\item Quelle partie du chemin matériel faudrait-il raccourcir ensuite, et que reste-t-il une fois la décision retirée~?
\item Que laisse de côté le modèle de course~?
\end{enumerate}

\medskip\noindent\textbf{Partie III --- L'emballement.}
\begin{enumerate}[resume]
\item Un moniteur logiciel met \qty{20}{\micro\second} à repérer un emballement et à positionner le bit kill. Combien d'ordres peuvent sortir entre-temps~?
\item Et avec un moniteur sur la carte elle-même, qui réagit en 16 cycles~?
\item Pourquoi le bit kill agit-il sur l'étape 3 plutôt que seulement sur l'étape 2~?
\item Quelles limites faudrait-il ajouter à la carte, et lesquelles doivent rester dans la couche du courtier~?
\end{enumerate}

\medskip\noindent\textbf{Partie IV --- Le verdict.}
\begin{enumerate}[resume]
\item Énoncer le \emph{résultat nommé}~: les latences des deux chemins, leur probabilité de gagner face à un concurrent de médiane
  \qty{1}{\micro\second}, et les ordres que le seau laisse sortir avant un kill à \qty{20}{\micro\second}.
\item Le chemin de \qty{0.82}{\micro\second} vaut-il la peine d'être construit pour une stratégie qui affronte des concurrents à \qty{5}{\micro\second}~?
\item À quoi la firme renonce-t-elle en transférant la décision dans le matériel~?
\item À quelle fréquence le seuil peut-il changer, et qui le modifie~?
\item Pourquoi le logiciel doit-il recevoir une copie de chaque ordre que la carte envoie~?
\item Comment tester le circuit sur une journée enregistrée avant sa mise en production~?
\item Que provoquerait une panne du chemin d'écriture de la \omterm{def:nw:programmable-hardware-ii-trading-designs:regmap}{carte des registres}, et comment le circuit la détecterait-il~?
\item En une phrase~: que décide le matériel, et que décide le logiciel~?
\end{enumerate}
\end{problem}

\section{Questions d'entretien}

\begin{interviewq}[$\star$ \iqroles{développeur}]\label{iq:nw:programmable-hardware-ii-trading-designs:1}
Comment un FPGA analyse-t-il un flux de données de marché pendant que le paquet arrive encore~?
\end{interviewq}

\begin{interviewq}[$\star\star$ \iqroles{développeur}]\label{iq:nw:programmable-hardware-ii-trading-designs:2}
Quels contrôles de risque mettriez-vous dans le matériel, et comment prouvez-vous qu'ils ne peuvent pas être contournés~?
\end{interviewq}

\begin{interviewq}[$\star\star$ \iqroles{développeur, trader}]\label{iq:nw:programmable-hardware-ii-trading-designs:3}
Qu'est-ce qui reste dans le logiciel d'un \omterm{def:nw:programmable-hardware-ii-trading-designs:hybrid}{système de trading hybride}, et pourquoi~?
\end{interviewq}

\begin{interviewq}[$\star\star$ \iqroles{développeur}]\label{iq:nw:programmable-hardware-ii-trading-designs:4}
Comment envoyez-vous un ordre par TCP depuis un FPGA~?
\end{interviewq}

\begin{interviewq}[$\star\star\star$ \iqroles{développeur}]\label{iq:nw:programmable-hardware-ii-trading-designs:5}
Comment vérifieriez-vous que le \omterm{def:nw:programmable-hardware-ii-trading-designs:trigger}{déclencheur matériel} et son modèle logiciel concordent, et que feriez-vous lorsqu'ils divergent~?
\end{interviewq}

\begin{interviewq}[$\star\star$ \iqroles{développeur, trader}]\label{iq:nw:programmable-hardware-ii-trading-designs:6}
Un fournisseur annonce 14~nanosecondes de tick-to-trade. Quelles questions posez-vous~?
\end{interviewq}
